\chapter{Hardware Overview}
\label{ch:hardware}

This chapter describes the board at a block level. The complete hardware
design --- schematics, PCB, Gerbers, and bill of materials --- is
published in the repository under the CERN-OHL-S v2 licence.

\section{System block diagram}

DigiRadio is built around four devices coordinated by the ESP32-S3. The
tuner produces the audio stream, the DSP conditions and mixes it, and the
Bluetooth module transmits it; the host manages all three and provides
the network interface.

\begin{table}[htbp]
  \centering
  \begin{tabular}{@{}lll@{}}
    \drhead Device & Role & Interface \\
    \midrule
    Si4684            & DAB+/FM tuner              & SPI / I\textsuperscript{2}C \\
    ADAU1701          & Audio DSP (EQ + mixer)     & I\textsuperscript{2}C \\
    ESP32-S3-WROOM-1  & Host controller, Wi-Fi/BLE & ---            \\
    FSC-BT1035         & Bluetooth 5.2 (aptX Adaptive) & UART (AT)  \\
    \bottomrule
  \end{tabular}
  \caption{Principal integrated circuits and their roles.}
  \label{tab:hw-ics}
\end{table}

\section{Power}

The board is powered from USB-C. An AP63203 buck converter generates the
main rails from the USB input, feeding the digital and analogue sections.
The tuner's sensitive supplies are derived and filtered separately to
keep RF performance clean.

\section{RF and antennas}

Two 2.4\,GHz radios are present --- the ESP32-S3 and the Bluetooth module.
Their antennas are placed on opposite diagonal corners of the board to
maximise separation and minimise mutual desense. Each module keeps its
antenna keep-out clear of copper on all layers.

\begin{drnote}[Board]
The PCB is a six-layer stack-up with controlled impedance on the USB and
RF sections. The full fabrication data is in the repository.
\end{drnote}

\section{User interface and connectors}

Configuration is done over Wi-Fi through a built-in web interface (see
Chapter~\ref{ch:firmware}); no display is required on the board itself.
Audio reaches the listener wirelessly over Bluetooth.
